% SPDX-License-Identifier: AGPL-3.0-or-later
% Copyright (C) 2026 Sanjeeva Wijeyesakere
%
% XPreview+ User Manual.
% Build:  tectonic xpreview-plus-manual.tex      (or: latexmk -xelatex / -pdf)
\documentclass[11pt,a4paper]{article}

% ---------------------------------------------------------------- engine/fonts
\usepackage{iftex}
\ifXeTeX
  \usepackage{fontspec}
  \setmainfont{texgyrepagella}[Extension=.otf, UprightFont=*-regular, BoldFont=*-bold,
                               ItalicFont=*-italic, BoldItalicFont=*-bolditalic]
  \setsansfont{texgyreheros}[Extension=.otf, UprightFont=*-regular, BoldFont=*-bold,
                             ItalicFont=*-italic, BoldItalicFont=*-bolditalic, Scale=0.92]
  \setmonofont{texgyrecursor}[Extension=.otf, UprightFont=*-regular, BoldFont=*-bold, Scale=0.95]
\else
  \usepackage[T1]{fontenc}
  \usepackage[utf8]{inputenc}
  \usepackage{tgpagella}
  \usepackage[scale=0.92]{tgheros}
  \usepackage{tgcursor}
\fi

% ---------------------------------------------------------------- packages
\usepackage[margin=2.4cm,headheight=14pt]{geometry}
\usepackage{graphicx}
\usepackage[table]{xcolor}
\usepackage{booktabs,tabularx,longtable,array}
\usepackage{enumitem}
\usepackage{microtype}
\usepackage{fancyhdr}
\usepackage[most]{tcolorbox}
\usepackage{titlesec}
\usepackage[section]{placeins}  % keep figures inside their section
\usepackage[hidelinks]{hyperref}
\usepackage{bookmark}

\definecolor{brand}{HTML}{1F5FBF}
\definecolor{brandlight}{HTML}{E8F0FB}
\definecolor{notebg}{HTML}{FFF6D6}
\definecolor{noteframe}{HTML}{E2C766}
\definecolor{keybg}{HTML}{F4F4F6}
\definecolor{keyframe}{HTML}{B9BBC2}
\definecolor{rule}{HTML}{C9CED8}

\hypersetup{colorlinks=true, linkcolor=brand, urlcolor=brand,
  pdftitle={XPreview+ User Manual}, pdfauthor={Sanjeeva Wijeyesakere},
  pdfsubject={User manual for XPreview+ 1.0}}

\setlist{itemsep=2pt, topsep=4pt}
\setlength{\parindent}{0pt}
\setlength{\parskip}{6pt}
\renewcommand{\arraystretch}{1.25}

\titleformat{\section}{\Large\sffamily\bfseries\color{brand}}{\thesection}{0.8em}{}
\titleformat{\subsection}{\large\sffamily\bfseries}{\thesubsection}{0.8em}{}
\titleformat{\subsubsection}{\normalsize\sffamily\bfseries}{}{0em}{}
\titlespacing*{\section}{0pt}{18pt}{8pt}
\titlespacing*{\subsection}{0pt}{12pt}{4pt}
\titlespacing*{\subsubsection}{0pt}{10pt}{2pt}
\setcounter{tocdepth}{2}
\usepackage{tocloft}  % compact table of contents (fits one page)
\setlength{\cftbeforesecskip}{3.5pt}
\renewcommand{\cftsecfont}{\sffamily\bfseries}
\renewcommand{\cftsecpagefont}{\sffamily\bfseries}
\setlength{\cftbeforetoctitleskip}{0pt}
\renewcommand{\cfttoctitlefont}{\Large\sffamily\bfseries\color{brand}}

\pagestyle{fancy}
\fancyhf{}
\fancyhead[L]{\small\sffamily XPreview+ 1.0 User Manual}
\fancyhead[R]{\small\sffamily\nouppercase{\leftmark}}
\fancyfoot[C]{\small\sffamily\thepage}
\renewcommand{\headrulewidth}{0.4pt}
\renewcommand{\sectionmark}[1]{\markboth{#1}{}}

% ---------------------------------------------------------------- macros
\newcommand{\app}{XPreview+}
% Keyboard key cap, e.g. \key{Cmd}+\key{O}
\newcommand{\key}[1]{{\setlength{\fboxsep}{1.4pt}\setlength{\fboxrule}{0.5pt}%
  \raisebox{-0.1ex}{\fcolorbox{keyframe}{keybg}{\footnotesize\sffamily\strut #1}}}}
% Menu path, e.g. \menu{File > Export as PDF...}
\newcommand{\menu}[1]{\textsf{\textbf{#1}}}
\newcommand{\ui}[1]{\textsf{#1}}
\newcommand{\file}[1]{\texttt{#1}}

\newtcolorbox{note}[1][Note]{enhanced, breakable, colback=notebg, colframe=noteframe,
  boxrule=0.5pt, arc=2pt, left=6pt, right=6pt, top=4pt, bottom=4pt,
  fonttitle=\sffamily\bfseries\small, coltitle=black, attach title to upper={\ }, title={#1:}}
\newtcolorbox{tip}{enhanced, breakable, colback=brandlight, colframe=brand!45, boxrule=0.5pt,
  arc=2pt, left=6pt, right=6pt, top=4pt, bottom=4pt,
  fonttitle=\sffamily\bfseries\small, coltitle=black, attach title to upper={\ }, title={Tip:}}

\newcommand{\screenshot}[3]{% file, width, caption
  \begin{figure}[htbp]\centering
    % The frame adds 2*(\fboxsep+\fboxrule); keep the total within #2.
    \fcolorbox{rule}{white}{\includegraphics[width=\dimexpr#2-2\fboxsep-2\fboxrule\relax]{figures/#1}}
    \caption{#3}\label{fig:#1}
  \end{figure}}

\begin{document}

% ================================================================ title page
\begin{titlepage}
  \centering
  \vspace*{2.5cm}
  \includegraphics[width=4.2cm]{figures/icon.png}\par
  \vspace{1.2cm}
  {\sffamily\bfseries\fontsize{38}{44}\selectfont \app\par}
  \vspace{0.4cm}
  {\sffamily\Large User Manual\par}
  \vspace{0.3cm}
  {\sffamily\large Version 1.0\par}
  \vspace{1.6cm}
  {\large A fast viewer for PDF documents, images, RTF and OpenDocument text,\\
  with export of anything to PDF.\par}
  \vfill
  {\sffamily macOS \quad·\quad Windows \quad·\quad Linux\par}
  \vspace{0.6cm}
  {\small Copyright \textcopyright{} 2026 Sanjeeva Wijeyesakere.\\
  \app{} is free software under the GNU Affero General Public License, version 3 or later.\\
  \url{https://github.com/swijeyes/XPreview-Plus}\par}
\end{titlepage}

\tableofcontents
\clearpage

% ================================================================
\section{Introduction}

\app{} is a document and image viewer in the spirit of macOS Preview. It opens a
file almost instantly, shows it page by page, lets you search, select and copy
text, and can turn any file it opens into a PDF. It runs on macOS and Windows, and can be built for Linux.

\subsection{What \app{} can open}
\begin{itemize}
  \item \textbf{PDF documents}, including password-protected ones.
  \item \textbf{Images}: PNG, JPEG, GIF, WebP, TIFF (including multi-page TIFF), BMP,
        SVG and more; HEIC on macOS. See Section~\ref{sec:formats} for the complete list.
  \item \textbf{Rich Text Format} (\file{.rtf}) documents, such as those from TextEdit, WordPad or Word.
  \item \textbf{OpenDocument Text} (\file{.odt}, \file{.ott}, \file{.fodt}) documents,
        such as those from LibreOffice, Google Docs or Microsoft Word.
\end{itemize}

\subsection{Highlights}
\begin{itemize}
  \item \textbf{Page thumbnails} and a \textbf{table of contents} in the sidebar.
  \item \textbf{Export as PDF} for every file type.
  \item \textbf{Search} with highlighted results, \textbf{text selection} and \textbf{copy}.
  \item Smooth \textbf{zooming} (including trackpad pinch), \textbf{fit to width or page}, and \textbf{rotation}.
  \item \textbf{Tabs}, so several documents can be open in one window.
  \item \textbf{Private and safe}: \app{} never connects to the internet on its own, never runs
        macros or scripts, and treats every file as potentially hostile.
\end{itemize}

\screenshot{window-odt.png}{\linewidth}{\app{} showing an OpenDocument text file, with page
thumbnails in the sidebar.}

% ================================================================
\section{Installing \app{}}

\subsection{macOS}
\textbf{Requirements:} a Mac with Apple silicon (M-series) and a recent version of macOS.

\begin{enumerate}
  \item Download \file{XPreviewPlus-1.0-macos-arm64.dmg} from the project's
        \href{https://github.com/swijeyes/XPreview-Plus/releases/latest}{Releases page}.
  \item Double-click the downloaded file to open it.
  \item Drag the \ui{\app} icon onto the \ui{Applications} folder shown next to it.
  \item Eject the disk image (drag it to the Trash, or click the eject button in Finder).
\end{enumerate}

\subsubsection{Opening it the first time}
If the copy you downloaded is not notarised by Apple, macOS blocks it the first time
you open it and shows a message that it ``cannot verify'' the app. To allow it:
\begin{enumerate}
  \item Try to open \app{} once and close the warning.
  \item Open \menu{System Settings > Privacy \& Security}.
  \item Scroll down to the message about \app{} and click \ui{Open Anyway}.
  \item Confirm with your password or Touch ID.
\end{enumerate}
You only need to do this once.

\subsubsection{Removing \app{}}
Drag \ui{\app} from \ui{Applications} to the Trash. To also remove its settings,
see Section~\ref{sec:settings}.

\subsection{Windows}
\textbf{Requirements:} 64-bit Windows 10 or Windows 11.

\app{} for Windows is \emph{portable}. There is no installer, and it never writes to the
Windows registry.
\begin{enumerate}
  \item Download \file{XPreviewPlus-1.0-windows-x64-portable.zip} from the
        \href{https://github.com/swijeyes/XPreview-Plus/releases/latest}{Releases page}.
  \item Right-click the zip file, choose \ui{Extract All\ldots}, and pick a folder,
        for example your \file{Documents} folder or a USB stick.
  \item Open the extracted \file{XPreview+} folder and double-click \file{XPreview+.exe}.
\end{enumerate}

If Microsoft Defender SmartScreen shows ``Windows protected your PC'', click
\ui{More info} and then \ui{Run anyway}. This appears for new applications that are not yet
widely downloaded.

\begin{note}
Keep all the files in the \file{XPreview+} folder together; the program needs them.
Your settings are stored in a \file{settings} folder inside it. To remove \app{}, delete the folder.
\end{note}

\subsection{Linux}
There is no ready-made Linux package yet. \app{} can be built from source; see
the \file{README} in the project repository.

\subsection{Checking your download (optional)}
Each download on the Releases page has a matching \file{.sha256} file containing its
checksum. To check a download, open a terminal in the folder where you saved both files and run
\texttt{shasum -a 256 -c} followed by the name of the \file{.sha256} file (on Windows,
\texttt{CertUtil -hashfile \textit{file} SHA256}), and compare.

% ================================================================
\section{The \app{} window}\label{sec:window}

\screenshot{window-pdf.png}{\linewidth}{The main window with a four-page PDF. Thumbnails are on
the left, the page view on the right, and the status bar along the bottom.}

\subsection{Toolbar}
From left to right:
\begin{description}[style=nextline, leftmargin=2.2em, font=\sffamily\bfseries]
  \item[Sidebar button] Click to show or hide the sidebar. Click the small arrow next to it
        to choose between \ui{Thumbnails} and \ui{Table of Contents}.
  \item[Zoom out / Zoom in] Step the zoom down or up.
  \item[Rotate left / Rotate right] Turn the pages a quarter turn.
  \item[Export as PDF] Save the current document as a PDF (Section~\ref{sec:export}).
  \item[Search field] Find text in the document (Section~\ref{sec:search}). The number of
        results appears to its left.
\end{description}

\subsection{Sidebar}
The sidebar shows either \textbf{page thumbnails} or the \textbf{table of contents}. Drag the
thin divider between the sidebar and the page to make it wider or narrower; thumbnails grow
and shrink with it.

\subsection{Page view}
The main area shows the pages of your document one after another on a grey background. Scroll to move
through the document.

\subsection{Status bar}
The left side shows the current page, for example ``Page 2 of 9''. For single-page files it shows the file type instead.
The right side has the \textbf{zoom box}. It shows the current zoom, and you can
choose \ui{Fit Width}, \ui{Fit Page} or a percentage from it, or type a value such as
\texttt{140\%} and press \key{Return}.

\subsection{Information bar}
For RTF and ODT documents that \app{} had to simplify, a yellow bar appears above the page
(Section~\ref{sec:simplified}). Click the \ui{×} button on its right to dismiss it.

\subsection{Tabs}
Each file you open gets its own tab. The tab bar appears as soon as two or more files are open.
Click a tab to switch, drag tabs to reorder them, and click a tab's close button
(or press \key{Cmd}+\key{W}) to close it.

% ================================================================
\section{Opening files}

There are several ways to open a file:
\begin{itemize}
  \item Choose \menu{File > Open\ldots} (\key{Cmd}+\key{O} on macOS, \key{Ctrl}+\key{O} on
        Windows) and select one or more files.
  \item \textbf{Drag} files onto the \app{} window.
  \item On macOS, drag files onto the \app{} icon in the \textbf{Dock}.
  \item In Finder, right-click a file and choose \menu{Open With > \app}.
  \item Choose a recently opened file from \menu{File > Open Recent}. The list keeps the last
        twelve files; \ui{Clear Menu} empties it.
\end{itemize}

If a file is already open, \app{} switches to its tab instead of opening it twice.

\begin{tip}
To have a type of file always open in \app{} on macOS, select such a file in Finder, choose
\menu{File > Get Info}, pick \app{} under \ui{Open with}, and click \ui{Change All\ldots}.
\end{tip}

\subsection{Password-protected PDFs}
When a PDF is encrypted, \app{} asks for its password. If you cancel, the tab shows a message
instead of the document; close it and open the file again to retry.

\subsection{How \app{} recognises files}
\app{} looks at a file's \emph{contents}, not just its name, to decide what it is. If a file's
extension is wrong (for example an RTF document named \file{report.pdf}), \app{} still opens it
correctly and notes ``detected as \ldots'' in the status bar.

\subsection{Files that cannot be opened}
If a file is damaged, of an unsupported type, or too large to open safely, its tab shows a short explanation.
OpenDocument spreadsheets (\file{.ods}) and presentations
(\file{.odp}) are recognised but not supported in this version.

% ================================================================
\section{Viewing and navigating}

\subsection{Moving through pages}
\begin{itemize}
  \item Scroll with the trackpad, mouse wheel or scroll bar.
  \item Click a \textbf{thumbnail} to jump to that page. The highlighted thumbnail follows
        the page you are reading.
  \item Click an entry in the \textbf{table of contents}.
  \item Use the \menu{Go} menu: \ui{Previous Page}, \ui{Next Page}, \ui{First Page},
        \ui{Last Page} and \ui{Go to Page\ldots}.
  \item On the keyboard: \key{Space} and \key{Shift}+\key{Space} move a screen down or up;
        \key{Page Up}/\key{Page Down} and the arrow keys scroll; \key{Home} and \key{End} go to the
        first and last page. When the page fits the window's width, \key{←} and \key{→} move to
        the previous and next page.
\end{itemize}

\subsection{Table of contents}
For PDFs, the table of contents comes from the document's bookmarks. For RTF and ODT
documents, \app{} builds it from the headings. If a document has neither, the sidebar shows
``No table of contents''.

\subsection{Zooming}
\begin{itemize}
  \item Use the zoom buttons in the toolbar, or \menu{View > Zoom In} / \menu{Zoom Out}.
  \item \menu{View > Actual Size} shows the page at 100\%.
  \item \menu{View > Zoom to Fit} fits the whole page in the window; \menu{View > Fit Width}
        fits the page width. Both keep fitting as you resize the window.
  \item On a trackpad, \textbf{pinch} to zoom. Double-tap with two fingers to switch between fit
        width and a closer view.
  \item Hold \key{Cmd} (macOS) or \key{Ctrl} (Windows) and scroll to zoom around the pointer.
  \item Use the zoom box in the status bar.
\end{itemize}
Zoom ranges from 5\% to 3200\%. Documents open fitted to the window's width. Images open
fitted to the window, but are never enlarged beyond 100\%.

\subsection{Rotating}
\menu{Tools > Rotate Left} and \menu{Rotate Right} (\key{Cmd}+\key{L}, \key{Cmd}+\key{R}) turn
all pages. Rotation only affects how the document is shown; the original file is never changed.
Exported PDFs and printouts use the rotation you see.

\subsection{Viewing images}
\screenshot{window-image.png}{0.8\linewidth}{An SVG image. Images are a single page and are
shown at their real size if they fit.}

When an image is larger than the window, \textbf{drag} it to move it around. \textbf{Double-click}
to switch between fitting the window and actual size. Multi-page TIFF files show
one page per image. Animated GIF and WebP images show their first frame.

% ================================================================
\section{Searching}\label{sec:search}

\begin{enumerate}
  \item Click the search field in the toolbar, or press \key{Cmd}+\key{F} (\key{Ctrl}+\key{F}).
  \item Type the text you are looking for and press \key{Return}.
\end{enumerate}
Every match is highlighted in yellow and the current one in orange. The counter next to the
field shows, for example, ``3 of 12'', or ``Not found''.

\begin{itemize}
  \item Press \key{Return} again, or choose \menu{Edit > Find Next}, to go to the next match.
  \item \key{Shift}+\key{Return} or \menu{Edit > Find Previous} goes back.
  \item Searching ignores upper and lower case.
  \item Clear the field (click its \ui{×}) to remove the highlights.
\end{itemize}

\begin{note}
Images have no text to search. Scanned PDFs that contain only pictures of pages cannot be
searched either, because \app{} does not perform text recognition (OCR).
\end{note}

% ================================================================
\section{Selecting and copying}

\subsection{Text}
In PDF, RTF and ODT documents, drag across text to select it; the selection is shown in blue.
Then choose \menu{Edit > Copy} (\key{Cmd}+\key{C}) to copy it. A selection covers text on one
page at a time. Click anywhere or press \key{Esc} to clear it.

\begin{tip}
To move around a zoomed-in document instead of selecting text, hold \key{Option} (\key{Alt} on
Windows) while you drag.
\end{tip}

\subsection{Images}
When an image is open and nothing is selected, \menu{Edit > Copy} copies the image itself, ready to paste
into another application.

\subsection{Links}
Click a link to follow it. Web and e-mail links (\texttt{http}, \texttt{https} and
\texttt{mailto}) open in your browser or mail app. Links inside a PDF jump to the page they
point to. The pointer turns into a hand over a link. For safety, other kinds of links are
ignored, and nothing is ever opened without a click.

% ================================================================
\section{Exporting as PDF}\label{sec:export}

Any file that \app{} can open can be saved as a PDF.
\begin{enumerate}
  \item Choose \menu{File > Export as PDF\ldots} (\key{Shift}+\key{Cmd}+\key{E}), or click the
        export button in the toolbar.
  \item Choose a name and location. \app{} suggests the original name with \file{.pdf}, in the same folder.
  \item Click \ui{Save}.
\end{enumerate}
A progress window appears for longer documents, and the status bar confirms when the export
is done.

\subsection{What you get}
\begin{description}[style=nextline, leftmargin=2.2em, font=\sffamily\bfseries]
  \item[PDF documents] An exact copy of the original file. If you rotated the view,
        the copy is saved with that rotation.
  \item[RTF and ODT documents] Real PDF pages with sharp, selectable and searchable text, the
        same layout you see in \app{}, and the fonts embedded.
  \item[Images] A PDF page the size of the image, with the picture at full quality. SVG images
        stay sharp at any zoom.
\end{description}

\begin{note}
Exports are written safely: the new PDF only replaces an existing file once it has been
written completely, so a failed or interrupted export never leaves a broken file behind.
\end{note}

% ================================================================
\section{Printing}

Choose \menu{File > Print\ldots} (\key{Cmd}+\key{P}). In the print dialog you can pick a printer
and the range of pages. Each page is scaled to fit the paper and centred, using the current
rotation.

% ================================================================
\section{Document information and the ``simplified view''}\label{sec:simplified}

\subsection{The Inspector}
\menu{Tools > Show Inspector} (\key{Cmd}+\key{I}) shows information about the open file: its type,
the number of pages and details such as title, author and dates. For images it shows the
dimensions and colour profile. It also lists anything that was shown in a simplified way.

\subsection{Simplified view}
RTF and ODT files are word-processor documents. \app{} contains its own renderer for them, so
you do not need LibreOffice or Microsoft Word installed. It reproduces text, fonts, colours, lists,
tables, images, headers and footers faithfully, but a few features are approximated. When
that happens, a yellow bar says
\emph{``Simplified view: some formatting may differ from the original.''} Its buttons are:
\begin{description}[style=nextline, leftmargin=2.2em, font=\sffamily\bfseries]
  \item[Details\ldots] Lists exactly what was simplified, for example ``Footnotes are shown at
        the end of the document'' or ``Font substituted: Calibri → Carlito''.
  \item[Open in Default App] Opens the file in the application your computer normally uses for it.
  \item[Export as PDF\ldots] Saves the document as you see it.
\end{description}
\menu{File > Open in Default App} is also available for any file.

\subsection{Fonts}
Documents often use fonts that are not installed on your computer. \app{} includes free fonts
with the same letter widths as the common Microsoft fonts, so lines and pages break in the same
places as in the original:

\begin{center}
\begin{tabular}{ll}
\toprule
\textbf{Font in the document} & \textbf{Shown with} \\
\midrule
Arial, Helvetica & Arimo \\
Times New Roman, Times & Tinos \\
Courier New, Courier & Cousine \\
Calibri & Carlito \\
Cambria & Caladea \\
\bottomrule
\end{tabular}
\end{center}
A font that is installed on your computer is always used as is.

% ================================================================
\section{Supported file types}\label{sec:formats}

\begin{longtable}{>{\raggedright\arraybackslash}p{3.3cm}>{\raggedright\arraybackslash}p{3.2cm}>{\raggedright\arraybackslash}p{8.2cm}}
\toprule
\textbf{Type} & \textbf{Extensions} & \textbf{Notes} \\
\midrule
\endhead
PDF & \file{.pdf} & Full support, including password-protected files, bookmarks and links. \\
PNG, JPEG, BMP & \file{.png .jpg .jpeg .bmp} & Rotation stored in photos (EXIF) is applied automatically. \\
GIF, WebP & \file{.gif .webp} & Animated images show their first frame. \\
TIFF & \file{.tif .tiff} & Multi-page TIFF files show all pages. \\
SVG & \file{.svg .svgz} & Vector images; stay sharp at every zoom. \\
HEIC, JPEG 2000 & \file{.heic .heif .jp2} & macOS only. \\
Other images & \file{.ico .icns .tga .wbmp} & \\
Rich Text Format & \file{.rtf} & Fonts, colours, lists, tables, pictures (PNG, JPEG), links, headers, footers, footnotes. \\
OpenDocument Text & \file{.odt .ott .fodt} & Styles, headings, lists, tables, images, links, headers, footers, page numbers, footnotes. \\
OpenDocument Spreadsheet / Presentation & \file{.ods .odp} & Recognised, but not supported yet. \\
\bottomrule
\end{longtable}

\subsection{Known limitations of the RTF and ODT renderer}
\begin{itemize}
  \item Multiple text columns are shown as a single column.
  \item Footnotes are collected at the end of the document rather than at the foot of each page.
  \item Text boxes and frames are placed in the flow of the text, near where they are anchored.
  \item Drawings made of shapes, embedded charts without a preview picture, and WMF/EMF pictures
        (common in older RTF files) appear as labelled placeholders.
  \item All pages use the page size of the document's first section.
  \item Comments are hidden. Documents with tracked changes show the final text.
\end{itemize}

% ================================================================
\section{Keyboard shortcuts}\label{sec:shortcuts}

On macOS, \key{Cmd} is the Command key and \key{Option} is the Option (Alt) key.

\begin{longtable}{>{\raggedright\arraybackslash}p{5.6cm}>{\raggedright\arraybackslash}p{4.6cm}>{\raggedright\arraybackslash}p{4.6cm}}
\toprule
\textbf{Action} & \textbf{macOS} & \textbf{Windows / Linux} \\
\midrule
\endhead
\multicolumn{3}{l}{\sffamily\bfseries File} \\
Open & \key{Cmd}+\key{O} & \key{Ctrl}+\key{O} \\
Close tab & \key{Cmd}+\key{W} & \key{Ctrl}+\key{W} \\
Export as PDF & \key{Shift}+\key{Cmd}+\key{E} & \key{Ctrl}+\key{Shift}+\key{E} \\
Print & \key{Cmd}+\key{P} & \key{Ctrl}+\key{P} \\
Quit & \key{Cmd}+\key{Q} & \key{Ctrl}+\key{Q} on Linux; close the window on Windows \\
\midrule
\multicolumn{3}{l}{\sffamily\bfseries Edit} \\
Copy & \key{Cmd}+\key{C} & \key{Ctrl}+\key{C} \\
Find & \key{Cmd}+\key{F} & \key{Ctrl}+\key{F} \\
Find next & \key{Cmd}+\key{G} & \key{F3} or \key{Ctrl}+\key{G} \\
Find previous & \key{Shift}+\key{Cmd}+\key{G} & \key{Shift}+\key{F3} or \key{Ctrl}+\key{Shift}+\key{G} \\
\midrule
\multicolumn{3}{l}{\sffamily\bfseries View} \\
Hide sidebar & \key{Option}+\key{Cmd}+\key{1} & \key{Ctrl}+\key{Alt}+\key{1} \\
Thumbnails & \key{Option}+\key{Cmd}+\key{2} & \key{Ctrl}+\key{Alt}+\key{2} \\
Table of contents & \key{Option}+\key{Cmd}+\key{3} & \key{Ctrl}+\key{Alt}+\key{3} \\
Actual size & \key{Cmd}+\key{0} & \key{Ctrl}+\key{0} \\
Zoom in & \key{Cmd}+\key{+} or \key{Cmd}+\key{=} & \key{Ctrl}+\key{+} \\
Zoom out & \key{Cmd}+\key{-} & \key{Ctrl}+\key{-} \\
Zoom to fit & \key{Cmd}+\key{9} & \key{Ctrl}+\key{9} \\
Fit width & \key{Cmd}+\key{8} & \key{Ctrl}+\key{8} \\
\midrule
\multicolumn{3}{l}{\sffamily\bfseries Go} \\
Previous / next page & \key{Option}+\key{Cmd}+\key{↑} / \key{↓} & \key{Ctrl}+\key{Alt}+\key{↑} / \key{↓} \\
First / last page & \key{Cmd}+\key{↑} / \key{Cmd}+\key{↓} & \key{Ctrl}+\key{↑} / \key{Ctrl}+\key{↓} \\
Go to page & \key{Option}+\key{Cmd}+\key{G} & \key{Ctrl}+\key{Alt}+\key{G} \\
\midrule
\multicolumn{3}{l}{\sffamily\bfseries Tools} \\
Rotate left / right & \key{Cmd}+\key{L} / \key{Cmd}+\key{R} & \key{Ctrl}+\key{L} / \key{Ctrl}+\key{R} \\
Show Inspector & \key{Cmd}+\key{I} & \key{Ctrl}+\key{I} \\
\midrule
\multicolumn{3}{l}{\sffamily\bfseries In the page view} \\
Screen down / up & \key{Space} / \key{Shift}+\key{Space} & \key{Space} / \key{Shift}+\key{Space} \\
First / last page & \key{Home} / \key{End} & \key{Home} / \key{End} \\
Previous / next page (page fits width) & \key{←} / \key{→} & \key{←} / \key{→} \\
Clear selection & \key{Esc} & \key{Esc} \\
Zoom around pointer & \key{Cmd}+scroll, or pinch & \key{Ctrl}+scroll \\
Move a zoomed document & \key{Option}+drag & \key{Alt}+drag \\
\bottomrule
\end{longtable}

% ================================================================
\section{Settings and privacy}\label{sec:settings}

\subsection{What \app{} remembers}
\app{} has no preferences window; it simply remembers:
\begin{itemize}
  \item the list of recently opened files,
  \item the sidebar mode (hidden, thumbnails or table of contents),
  \item the window's size and position,
  \item the folder you last opened a file from.
\end{itemize}
These are stored in:
\begin{description}[style=nextline, leftmargin=2.2em, font=\sffamily\bfseries]
  \item[macOS] \file{\textasciitilde/Library/Preferences/org.xpreview.XPreview.plist}
  \item[Windows (portable)] \file{settings\textbackslash XPreview\textbackslash XPreview.ini} inside the
        \file{XPreview+} folder
  \item[Linux] \file{\textasciitilde/.config/XPreview/XPreview.conf}
\end{description}
To reset \app{}, quit it and delete that file. On macOS you can instead run
\texttt{defaults delete org.xpreview.XPreview} in Terminal.

\subsection{Privacy and safety}
\begin{itemize}
  \item \app{} never connects to the internet by itself. It contains no analytics, no update
        checks and no accounts.
  \item Documents cannot make \app{} load anything from the internet or from elsewhere on
        your computer. For example, pictures referenced by an SVG or ODT file but stored
        outside it are not loaded.
  \item Macros and scripts in documents are never run.
  \item Links only open when you click them, and only web and e-mail links.
  \item Every file is treated as potentially hostile. Damaged or deliberately malformed files
        are handled safely and size limits stop them from exhausting memory. Such
        a file shows an error message instead of crashing the application.
  \item Exporting never changes the original file.
\end{itemize}

% ================================================================
\section{Troubleshooting}

\begin{description}[style=nextline, leftmargin=2.2em, font=\sffamily\bfseries]
  \item[macOS says \app{} cannot be opened or verified]
        Allow it once in \menu{System Settings > Privacy \& Security} (see Section~2.1).
  \item[Windows shows ``Windows protected your PC'']
        Click \ui{More info}, then \ui{Run anyway}.
  \item[\app{} is missing from ``Open With'']
        macOS adds \app{} to that list after it has been copied to \ui{Applications} and opened once.
  \item[A document looks different from the original]
        For RTF and ODT files, click \ui{Details\ldots} in the yellow bar to see what was
        simplified. For an exact rendering, use \ui{Open in Default App}.
  \item[Text cannot be found or selected]
        The document may be a scanned image without a text layer, or an image file.
  \item[``The file is too large to open'']
        \app{} limits very large or deeply nested files to stay safe and responsive. For
        example, images may be at most 30\,000 pixels wide or high.
  \item[A spreadsheet or presentation does not open]
        OpenDocument spreadsheets and presentations are not supported in version 1.0.
  \item[HEIC photos do not open on Windows]
        HEIC and JPEG 2000 images are only supported on macOS.
  \item[Something else is wrong]
        Reset \app{} as described in Section~\ref{sec:settings}. If the problem remains, please
        report it (below).
\end{description}

\subsection{Reporting problems}
Please report bugs and ask for features on the project's
\href{https://github.com/swijeyes/XPreview-Plus/issues}{GitHub issues page}. Include your
operating system, the \app{} version (\menu{Help > About \app}), and, if possible, a file that
shows the problem. Report \textbf{security} problems privately through the repository's
\ui{Security} tab instead.

% ================================================================
\section{License and credits}

\app{} is Copyright \textcopyright{} 2026 Sanjeeva Wijeyesakere.

\app{} is free software: you can redistribute it and/or modify it under the terms of the GNU
Affero General Public License as published by the Free Software Foundation, either version 3 of
the License, or (at your option) any later version. It is distributed in the hope that it will be
useful, but \textbf{without any warranty}; without even the implied warranty of merchantability or
fitness for a particular purpose. See the license for details. The full text is included with the
application (\file{LICENSE}). The source code is available at
\url{https://github.com/swijeyes/XPreview-Plus}.

\app{} is built with:
\begin{itemize}
  \item \textbf{Qt} (LGPL-3.0), the application framework;
  \item \textbf{PDFium} (BSD-3-Clause / Apache-2.0), for PDF rendering;
  \item \textbf{miniz} (MIT), for reading OpenDocument packages;
  \item the \textbf{Arimo}, \textbf{Tinos} and \textbf{Cousine} fonts (Apache-2.0), and the
        \textbf{Carlito} and \textbf{Caladea} fonts (SIL Open Font License 1.1).
\end{itemize}
Their licenses are listed in \file{THIRD\_PARTY\_NOTICES.md} and included with the application.
The \app{} icon was created by Sanjeeva Wijeyesakere with Google Gemini.

\end{document}
